\documentclass[10pt]{article}

\usepackage[utf8]{inputenc}

\usepackage[T1]{fontenc}

\usepackage{amsmath}

\usepackage{amsfonts}

\usepackage{amssymb}

\usepackage[version=4]{mhchem}

\usepackage{stmaryrd}

\usepackage{bbold}

\usepackage{mathrsfs}

\usepackage{graphicx}

\usepackage[export]{adjustbox}

\graphicspath{ {./images/} }



\DeclareUnicodeCharacter{0131}{$\imath$}



\begin{document}

щенная. Если $\sigma^{2}$ известна, $\bar{X}$ - несмещенная оценка минимальной дисперсии, $\bar{X}$ - минимаксная оценка по отношению к квадратичноӥ функции потерь.



2.2. Пусть $X_{j}$ - нормальные величины в $\mathbb{R}^{k}$ с плотностью



\[

(2 \pi)^{-k / 2} \exp \left\{-2^{-1}|x-\theta|^{2}\right\}, \theta \in \mathbb{R}^{k} .

\]



Статистика $\bar{X}$ - несмещенная оценка $\theta$, если $k \leqslant 2$ она допустима по отношению к квадратичной функции потерь, если $k>2$ - недопустима.



2.3. Пусть $X_{j}$ - случайные величины в $\mathbb{R} 1$ с неизвестной Іглотностью распределения $f$, принадлежащей заданному семейству $F$ плотностей. Для достаточно широкого класса $F$ это непараметритеская задача. Задача оценки значения плотности $f\left(x_{0}\right)$ в тотке $x_{0}-$ это задача оценки функционала $g(f)=f\left(x_{0}\right)$.



\begin{enumerate}

  \setcounter{enumi}{2}

  \item Лин ейная рег ре с с ионная модель. Наблюдаются величины

\end{enumerate}



\[

X_{i}=\sum_{\alpha=1}^{p} a_{\alpha i} \theta_{\alpha}+\xi_{i}

\]



$\xi_{i}$ - случайные возмущения, $i=\overline{1, n}$, матрица $\left\|a_{\alpha i}\right\|$ известна, оценке подлежит параметр $\left(\theta_{1}, \ldots, \theta_{p}\right)$.



\begin{enumerate}

  \setcounter{enumi}{3}

  \item Наблюдается отрезок гауссовского стационарного процесса $x(t), 0 \leqslant t \leqslant T$, с рациональной спектральной плотностью $\sum_{j=0}^{m} a_{j} \lambda j^{2} \cdot\left|\sum_{j=0}^{n} b_{j} \lambda^{j}\right|^{-2}$; оценке подлежат неизвестные параметры $\left\{a_{j}\right\},\left\{b_{j}\right\}$.

\end{enumerate}



С пос о б ы п о с т р о е ния оце нок. Наиболее распространенный максимального правдоподобия метод рекомендует в качестве оценки $\theta$ максимального правдоподобия оценку $\hat{\theta}(X)$, определяемую как точка максимума случайной функции $p(X ; \theta)$. Если $\Theta \subset \mathbb{R} k$, оценки максимального правдоподобия содержатся среди корнеӥ уравнения правдоподобия



\[

\frac{\partial}{\partial \theta} \ln p(\theta ; X)=0 .

\]



Наиженьших квадратов летод в применении к задаче примера 3 рекомендует в качестве оценки точку минимума функции



\[

m(\theta)=\sum_{1}^{n}\left(X_{1}-\sum_{\alpha} a_{\rho, \ell} \theta_{\alpha}\right)^{2} .

\]



Еще один способ построения состоит в том, что в качестве оценки для $\theta$ выбирается бейесовская по отношению к нек-рой функции потерь $w$ и нек-рому априорному распределению $Q$ оценка $T$, хотя исходная постановка не является бейесовской. Напр., если $\Theta=\mathbb{R}^{1}$ можно оценить $\theta$ посредством



\[

\int_{-\infty}^{\infty} 0_{p}(X ; \theta) d \theta\left(\int_{-\infty}^{\infty} p(X ; \theta) d \theta\right)^{-1} .

\]



Это бейесовская оценка по отношению к квадратической функции готерь и равномерному априорному распределению.



Маментов метод заключается в следующем. Пусть $\Theta \subset \mathbb{R}^{k}$ п пусть имеется $k$ «хороших» оценок $a_{1}(X), \ldots$, $a_{k}(X)$ для $\alpha_{1}(\theta), \ldots, \quad \alpha_{k}(\theta)$. Оценки метода моментов суть решения системы $\alpha_{i}(\theta)=a_{i}$. Часто в качестве $a_{i}$ выбирают эмширич. моменты (см. пример 1).



Если наб́людается выборка $X_{1}, \ldots, X_{n}$, то (см. пример 1) в качестве оценки для $g(\mathscr{P})$ можно выбрать $g\left(\mathscr{P}_{n}^{*}\right)$. Если функция $g\left(\mathscr{P}_{n}^{*}\right)$ не определена (напр., $g(\mathscr{P})=\frac{d \mathscr{P}}{d \lambda}(x), \lambda$ - мера ЈГебега), выбирают подхоцящие модификации $g_{n}\left(\mathscr{P}_{n}^{*}\right)$. Напр., дія оценки плотности употребляется гистогралља шли оценки вида



\[

\int \varphi_{n}(x-y) d \mathscr{P}_{n}^{*}(y) .

\]



А симптот ическое поведение оцен о к. Пусть для определенности рассматривается задача примера $2, \quad \Theta \subset \mathbb{R}^{k}$. Можно ожидать, что при $n \rightarrow \infty$ «хорошие» опенки должны неограниченно сближаться с оцениваемой характеристикой. Последовательность оценок $\theta_{n}^{*}\left(X_{1}, \ldots, X_{n}\right)$ наз. с о с т о яте льно й по сле д о в а те льнос т ь ю н о к $\theta$, если $\theta_{n}^{*} \rightarrow \theta$ по $P_{\theta}-$ вероятности для всех $\theta$. Перечисленные выше способы построения оценок в широких предположениях приводят к состоятельны.м оценкам. Оценки примера 1 состоятельны. Для регулярных задач оценивания оценки максимального правдоподобия и бейесовские асимптотически нормальны со средним $\theta$ и корреляционной матрицей $(n I(\theta))^{-1}$. В этих же условиях әти оценки асимптотически локально минимаксны по отнопению к широкому классу функций потерь и их можно рассматривать как асимптотически оптимальные (см. Асимптотически әффективная оценка).



Ин те р в а ль о е оцени в ани е. Случайное подмножество $E=E(X)$ множества $\Theta$ наз. д о в ери тельным множ е с т в м для оценки 0 с дов е рительным коэффициентом $ү$, если $P_{\theta}\{E \supset \theta\}=\gamma(\geqslant \vartheta)$. Обычно существует много доверительных областей с заданным $\gamma$ и проблема заключается в том, чтобы выбрать ту, к-рая обладает нек-рыми оптимальными свойствами (напр., интервал минимальної длины, если $\Theta \subset \mathbb{R}^{1}$ ). Пусть в условиях примера $2.1 \sigma=1$, тогда интервал



\[

\begin{gathered}

{\left[\bar{X}-\frac{\lambda}{\sqrt{n}}, \bar{X}+\frac{\lambda}{\sqrt{n}}\right],} \\

\gamma=\sqrt{\frac{2}{\pi}} \int_{\lambda}^{\infty} \exp \left\{-\frac{u^{2}}{2}\right\} d u

\end{gathered}

\]



является доверительным интервалом с доверительным коәффициентом $\gamma$ (см. Интервальная оуенка).



Jum.: [1] F i s h e r R., "Phil. Trans. Roy. Soc." (A), 1922,

t. 222, s. 309-68; [2] K o J м o r o p o в A. H., "Изв.



\includegraphics[max width=\textwidth, center]{2023_07_09_376efbb61a99bb50c6d7g-1}

Математические методы статистики, пер. с англ.,-М., 1948 ; [4] К е н д а лл М., С т ь ю а и связи, пер. с айГI., М., 1973; [5] Иб р а г и м о в И. А., ния, М., 1979; [6] Ч е н ц о в Н. Н, Статистические решающие правила и оптимальные выводы, М., 1972; [7] 3 а к с III., Теория статистических выводов, пер. с англ., M., 1975; [8] Gr e n-



СТАТИСТИЧЕСКОЙ МЕХАНИКИ МАТЕМАТИЧЕСКИЕ ЗАДАЧИ - совокупность общих проблем математич. физики, возникших из стремления четко осмыслить основные концепции и факты статистич. механики. Эти проблемы можно условно разделить на следующие группы:



\begin{enumerate}

  \item обоснование основных принципов статистич. механики,



  \item равновесные ансамбли в термодинамич. гределе, вывод термодинамич. соотношений,



  \item фазовые переходы,



  \item эволюция ансамблеї, проблема релаксации, исследование кинетич. и гидродинамич. уравнений,



  \item основные состояния, әлгментарные возбужкдения (в случае квантовых систем).



\end{enumerate}



Статистич. механика изучает системт, состоящие из большого числа (микроскопических) частиц, заключенных внутри большой (сравнительно с характерными размерами частиц) области $V$ ıространства $\mathbb{R}^{\mathbf{3}}$. Статистич. механика - в зависимости от способа описания системы - разделяется на классическую и квантовую.



Описание классич. системы, заключенной в области $V$, включает указание пространства $X$ возможных состояний каждой отдельної частицы (одночаститное пространство), а также совокупности $\Omega_{V} \cup N \geqslant 0 X^{N}$ допустимых конфигураций $\omega=\left\{x_{1}, \ldots, x_{N}\right\} ; x \in X N$;





\end{document}